\ncdlesson{7}{Cykel eller metro til arbejde?}{骑车还是坐地铁去上班？}
{DU3 Modul 1--2 交通话题｜对标 NCE 1 Lesson 33--38}
{五大情态动词（\textit{kan, skal, vil, må, bør} + 不带 \textit{at} 的动词原形）与中心状语夹心位置 · 交通方式固定介词搭配（\textit{med bus/tog/metro, på cykel, i bil, til fods}）}
{交通通勤是 DU3 模块一二最高频的口语话题！请牢记两条死规矩：第一，情态动词后面接原形时\textbf{绝对不能加 \textit{at}}；第二，\textit{ikke / altid} 夹在“情态动词 + 原形”中间。同时练熟情态动词在句中的极度弱读！}

% ---------------------------------------------------------------------
\begin{foerduleaeser}
带着下面两个问题阅读课文对话：
\begin{enumerate}
  \item \textit{Hvorfor vil de ikke cykle eller tage bussen til arbejde i dag?}（为什么他们今天不骑车也不坐公交车去上班？）
  \item \textit{Hvilket transportmiddel vælger de til sidst, og hvordan kommer de hen til stationen?}（他们最后选择了哪种交通工具？怎么去地铁站？）
\end{enumerate}
\end{foerduleaeser}

% ---------------------------------------------------------------------
\ncdsubsec{二、Dialog \& Udtale（课文对话、发音标注与中文译文）}

\begin{tipbox}{场景说明}
工作日清晨，窗外下着大雨还刮着风。Emma 和同事兼合租邻居 Jonas 站在玄关窗边，商量今天改用哪种交通方式去上班。
\end{tipbox}

\dialogturn{1}{Jonas}{Se ud ad vinduet, Emma! Det regner kraftigt, og vinden blæser. Vi kan ikke tage på arbejde på cykel i dag.}
{\textit{vinduet} \ipa{ˈvenduːˀəð}（词尾 \textit{-et} 读 \ipa{əð}）；\textit{regner} \ipa{ˈʁɑɪ̯nɐ}（\textit{eg} 读 \ipa{ɑɪ̯}）；\textit{kraftigt} \ipa{ˈkʰʁafdi}（\textbf{g 和 t 不发音}）；\textit{kan ikke} 在快速口语里常缩读为 \ipa{kʰa eg}（\textit{kan} 吞掉 \textit{n}）。}
{艾玛，看看窗外！雨下得好大，而且还在刮风。今天我们没法骑车去上班了。}

\dialogturn{2}{Emma}{Øv, jeg vil ellers helst cykle, for så får jeg motion. Men du har ret --- vi må hellere tage offentlig transport.}
{\textit{cykle} \ipa{ˈsyɡlə}（\textit{c} 读 \ipa{s}，\textit{y} 圆唇）；\textit{motion} \ipa{moˈɕoːˀn}（\textit{-ti-} 读 \ipa{ɕ}）；\textit{offentlig} \ipa{ˈʌfənli}（\textbf{t 和 g 均不发音}！）。}
{哎呀，要不然我其实最想骑车，因为那样能锻炼身体。不过你说得对——我们还是坐公共交通吧。}

\dialogturn{3}{Jonas}{Skal vi køre med metro, eller vil du hellere med bus eller tog?}
{\textit{skal vi} 在句中常弱读为 \ipa{sga vi}（\textbf{l 脱落}）；\textit{metro} \ipa{ˈmeːtʁo}（重音在第一音节）；\textit{tog} \ipa{ˈtʰɔʊ̯ˀ}（\textit{g} 滑化成 \ipa{ʊ̯}）。}
{咱们坐地铁去，还是你更想坐公交或者火车？}

\dialogturn{4}{Emma}{Vi bør ikke tage bussen i myldretiden, for den sidder ofte fast i trafikken.}
{\textit{bør} \ipa{ˈbɶɐ̯}（短促圆唇，带微卷收尾）；\textit{myldretiden} \ipa{ˈmylʁəˌtʰiːðˀən}（\textbf{d in \textit{myldre} 静音}）；\textit{trafikken} \ipa{tʰʁaˈfiɡən}（重音在第二音节）。}
{高峰期咱们不应该坐公交车，因为它经常堵在路上。}

\dialogturn{5}{Jonas}{Enig! Så tager vi metroen! Den kører hurtigt, og vi skal ikke skifte undervejs.}
{\textit{enig} \ipa{ˈeːni}（\textbf{词尾 g 静音}！）；\textit{metroen} \ipa{ˈmeːtʁoən}；\textit{skifte} \ipa{ˈsgifdə}；\textit{undervejs} \ipa{ʌnɐˈvɑɪ̯ˀs}（\textbf{d 不发音}）。}
{同意！那咱们坐地铁！地铁开得快，而且途中不用换乘。}

\dialogturn{6}{Emma}{Hvor langt er der til metrostationen? Kan vi gå til fods på ti minutter?}
{\textit{langt} \ipa{ˈlaŋˀd}；\textit{til fods} \ipa{tʰel ˈfoːˀs}（\textbf{d 不发音}，带喉塞音 \ipa{ˀ}）。}
{到地铁站有多远？我们步行十分钟能走到吗？}

\dialogturn{7}{Jonas}{Ja, vi kan godt gå til fods med en paraply. Det tager kun otte minutter.}
{\textit{kan godt} \ipa{kʰa ˈɡʌd}（注意 \textit{kan} 弱读，重音放在 \textit{godt} 上，\textit{godt} 读 \ipa{ɡʌd} 表示“完全可以”）；\textit{paraply} \ipa{pʰɑʁɑˈpʰlyːˀ}（重音在最后）。}
{可以，打把伞完全能走过去。只需要八分钟。}

\dialogturn{8}{Emma}{Hvad med min nabo Sara? Hun kører altid i bil til arbejde, men man må ikke parkere gratis derinde.}
{\textit{i bil} \ipa{i ˈbiːˀl}；\textit{må ikke} \ipa{mɔ ˈeɡə}；\textit{parkere} \ipa{pʰɑˈkʰeːˀɐ}；\textit{gratis} \ipa{ˈɡʁæːtis}（重音在第一音节）。}
{我邻居萨拉呢？她总是开车去上班，但市中心那边不准免费停车呀。}

\dialogturn{9}{Jonas}{Præcis! Derfor bør flere tage toget eller metroen, når vejret er dårligt. Nu må vi altså gå!}
{\textit{derfor} \ipa{ˈdɛɐ̯ˌfɒ}（注意重音在第一音节，触发 V2 倒装）；\textit{vejret} \ipa{ˈvɛɐ̯ˀð}（\textbf{j} 融于元音中，尾音软 \ipa{ð}）；\textit{altså} \ipa{ˈalˀsɔ}（\textbf{t 不发音}）。}
{没错！所以天气不好的时候，应该有更多人坐火车或地铁。好了，咱们现在真的得走了！}

\dialogturn{10}{Emma}{Ja, tag din jakke på, så løber vi ned til stationen!}
{\textit{tag ... på} \ipa{ˈtʰæː ... ˈpʰɔːˀ}（祈使句中 \textit{tag} 常丢掉 \textit{g}）；\textit{jakke} \ipa{ˈjaɡə}；\textit{løber} \ipa{ˈløːˀbɐ}。}
{好，穿上你的外套，咱们跑着去车站吧！}

% ---------------------------------------------------------------------
\ncdsubsec{三、Gloser（核心生词表）}

\begin{tipbox}{词汇提示}
五大情态动词的不定式（\textit{at kunne} 等）很少单独使用，请重点记熟括号第一项的\textbf{现在时词形（\textit{kan, skal, vil, må, bør}）}。注意 \textit{at cykle} 表到达某地时完成时用 \dabold{er}。
\end{tipbox}

\begin{glosertable}
情态动词 & at kunne & \textit{(kan, kunne, har kunnet)} & 能，会，可以（能力/可能） \\
情态动词 & at skulle & \textit{(skal, skulle, har skullet)} & 必须；打算，将要 \\
情态动词 & at ville & \textit{(vil, ville, har villet)} & 想要，愿意 \\
情态动词 & at måtte & \textit{(må, måtte, har måttet)} & 可以（允许）；必须（\textit{må ikke} 禁止） \\
情态动词 & at burde & \textit{(bør, burde, har burdet)} & 应该，应当（客观建议） \\
\midrule
动词 & at regne & \textit{(regner, regnede, har regnet)} & 下雨（\textit{Det regner.}） \\
动词 & at blæse & \textit{(blæser, blæste, har blæst)} & 刮风（\textit{Det blæser.}） \\
动词 \dabold{[er]} & at cykle & \textit{(cykler, cyklede, \dabold{er cyklet} / har cyklet)} & 骑自行车（到某地用 \textit{er}） \\
动词 & at skifte & \textit{(skifter, skiftede, har skiftet)} & 换乘；更换 \\
动词 & at parkere & \textit{(parkerer, parkerede, har parkeret)} & 停车，泊车 \\
\midrule
名词 & en cykel & \textit{(cyklen, cykler, cyklerne)} & 自行车（变形时省去 \textit{-e-}） \\
名词 & en metro & \textit{(metroen, metroer, metroerne)} & 地铁 \\
名词 & en bil & \textit{(bilen, biler, bilerne)} & 汽车，私家车 \\
名词 & et vejr & \textit{(vejret, 不可数)} & 天气 \\
名词 & en myldretid & \textit{(myldretiden, 不可数)} & 交通高峰期 \\
名词 & en paraply & \textit{(paraplyen, paraplyer, paraplyerne)} & 雨伞 \\
名词 & et vindue & \textit{(vinduet, vinduer, vinduerne)} & 窗户 \\
名词 & en motion & \textit{(motionen, 不可数)} & 运动，身体锻炼 \\
\end{glosertable}

% ---------------------------------------------------------------------
\ncdsubsec{四、Grammatik \& Faste vendinger（核心语法与地道表达精讲）}

\begin{grammatikbox}{语法点 1：五大情态动词（Modalverber）的意义与两大死命令}

\vspace{0.2em}
\renewcommand{\arraystretch}{1.24}
{\small
\begin{tabularx}{\linewidth}{@{} >{\itshape\bfseries\color{NCDNavy}}p{1.2cm} >{\raggedright\arraybackslash}p{4.3cm} >{\raggedright\arraybackslash}X @{}}
\toprule
\textbf{词形} & \textbf{核心语义} & \textbf{课文例句与考点辨析} \\
\midrule
kan & 有能力、条件允许 & \textit{Vi \da{kan} ikke \dabold{tage} på arbejde på cykel i dag.} \\
skal & 计划、安排、必然要求 & \textit{\da{Skal} vi \dabold{køre} med metro?} / \textit{Vi \da{skal} ikke \dabold{skifte}.} \\
vil & 个人主观意愿、想要 & \textit{Jeg \da{vil} ellers helst \dabold{cykle}.} \\
må & 允许、必要性；\dabold{må ikke = 禁止} & \textit{Man \da{må ikke} \dabold{parkere} gratis.}（不准免费停） \\
bør & 建议、理应、最好 & \textit{Vi \da{bør} ikke \dabold{tage} bussen i myldretiden.} \\
\bottomrule
\end{tabularx}
}
\vspace{0.25em}

\textbf{两大语法死命令（考试必考）：}
\begin{enumerate}
  \item \textbf{后接不带 \textit{at} 的动词原形！}：虽然普通动词后常带 \textit{at}（如 \textit{at lære}），但\textbf{五大情态动词后面绝对不加 \textit{at}}！\newline 对：\textit{Jeg kan \da{cykle}.} \quad 错：*\textit{Jeg kan at cykle.} / *\textit{Jeg kan cykler.}
  \item \textbf{中心状语夹在中间！}：\textit{ikke, altid, aldrig, ofte, godt, helst} 必须紧跟在\textbf{情态动词之后、主要动词原形之前}！\newline 陈述句：\textit{Vi \da{kan} \dabold{ikke} \da{cykle}.} \quad 倒装句：\textit{I dag \da{kan} vi \dabold{ikke} \da{cykle}.}
\end{enumerate}
\textbf{必考辨析}：英语“You must not”表禁止，在丹麦语中必须用 \dabold{må ikke}（禁止/不准）；如果说成 \textit{skal ikke}，意思只是“不必/不需要/不打算”。
\end{grammatikbox}

\begin{grammatikbox}{语法点 2：交通方式的四大固定介词搭配（绝不能混用！）}
在 DU3 口语和写作中，介绍怎么上班或上学，名词的冠词和介词搭配有一套固定组合：

\vspace{0.2em}
\renewcommand{\arraystretch}{1.24}
{\small
\begin{tabularx}{\linewidth}{@{} >{\bfseries\color{NCDRed}}p{2.3cm} >{\raggedright\arraybackslash}p{4.8cm} >{\raggedright\arraybackslash}X @{}}
\toprule
\textbf{介词} & \textbf{适用交通工具} & \textbf{固定例句（注意有无冠词后缀！）} \\
\midrule
\da{med} + 零冠词 & 公共交通（bus, tog, metro, fly） & \textit{Jeg kører \da{med bus / med tog / med metro}.} \\
\da{på} + 零冠词 & 跨骑/敞篷两轮代步工具 & \textit{Jeg tager på arbejde \da{på cykel}.}（= \textit{Jeg cykler.}） \\
\da{i} + 零冠词 & 封闭式私人汽车、出租车 & \textit{Hun kører altid \da{i bil} / \da{i taxa}.} \\
\da{til} + 特定形式 & 步行 & \textit{Vi går \da{til fods}.}（固定短语：步行） \\
\bottomrule
\end{tabularx}
}
\vspace{0.2em}

\textbf{一秒对比}：如果不用介词，而是直接用动词 \dabold{at tage}（乘坐），后面的公共交通工具\textbf{必须加定冠词后缀}：\textit{tage \dabold{bussen / toget / metroen}}（对比：\textit{køre med bus / tog / metro} 不加词尾！）。
\end{grammatikbox}

\begin{grammatikbox}{语法点 3：高频地道口语（Faste vendinger）}
\begin{itemize}
  \item \da{Det regner / Det blæser.} —— 下雨了 / 刮风了（天气形式主语固定用 \textit{det}）。
  \item \da{at have ret} / \da{Enig!} —— 说得对 / 同意！
  \item \da{at sidde fast i trafikken} —— 堵在车流里（堵车）。
  \item \da{at få motion} —— 得到锻炼（\textit{Jeg cykler, for så får jeg motion.}）。
  \item \da{kan godt / vil hellere / må hellere} —— “完全可以 / 更愿意 / 还是……比较好”（\textit{Vi må hellere tage metroen} 咱们还是坐地铁吧）。
\end{itemize}
\end{grammatikbox}

% ---------------------------------------------------------------------
\ncdsubsec{五、Mønsterøvelser（新概念式句型操练）}

\textbf{练习 A：用正确的情态动词（\textit{kan, skal, vil, må, bør}）与括号内动词的正确形式填空}
\begin{enumerate}
  \item Man \underline{\hspace{1.4cm}} \dabold{ikke} \underline{\hspace{2.0cm}} (\textit{at parkere}) her --- det er forbudt!（禁止）
  \item Vejret er dårligt, så vi \underline{\hspace{1.4cm}} hellere \underline{\hspace{2.0cm}} (\textit{at tage}) metroen.（建议/最好）
  \item Jeg \underline{\hspace{1.4cm}} godt \underline{\hspace{2.0cm}} (\textit{at tale}) lidt dansk nu.（能力）
  \item \underline{\hspace{1.4cm}} du \underline{\hspace{2.0cm}} (\textit{at køre}) med bus eller gå til fods?（意愿）
  \item Om mandagen \underline{\hspace{1.4cm}} Emma \underline{\hspace{2.0cm}} (\textit{at møde}) klokken halv ni.（计划/义务）
\end{enumerate}

\textbf{练习 B：交通介词与冠词实战选填（\textit{med, på, i, til} 或判断是否加 \textit{-en/-et}）}
\begin{enumerate}
  \item Om sommeren tager jeg næsten altid på arbejde \underline{\hspace{1.4cm}} cykel.
  \item Min bror bor langt væk, så han kører altid \underline{\hspace{1.4cm}} bil.
  \item Når det regner, tager jeg \underline{\hspace{2.4cm}} (\textit{tog / toget}) i stedet for at cykle.
  \item Kører du \underline{\hspace{1.4cm}} metro eller \underline{\hspace{1.4cm}} bus til sprogskolen?
  \item Metrostationen ligger lige her, så vi kan gå \underline{\hspace{1.4cm}} fods.
\end{enumerate}

\textbf{练习 C：中译丹实战（注意 V2 倒装、情态动词与 \textit{ikke} 位置）}
\begin{enumerate}
  \item 今天我们不能骑自行车去上班。（用 \textit{I dag} 开头触发倒装）
  \item 坐地铁我们不用换乘。（用 \textit{skal ikke}）
  \item 高峰期大家不应该开车去市中心。
  \item 你更想坐火车还是坐地铁？
  \item 在这里不准免费停车，所以很多人坐公交车。
\end{enumerate}

% ---------------------------------------------------------------------
\ncdsubsec{六、Resumé \& Mundtlig træning（口语复述与参考答案）}

\textbf{口语复述任务}：遮住课文，先口头回答 \textbf{Før du læser} 的两个问题，然后用 \textbf{4--5 句话} 分别描述你在晴天和雨天各自如何通勤（至少用到 3 个情态动词和 3 种不同的交通表达）。

\begin{answerbox}
\textbf{练习 A 答案}
\begin{enumerate}
  \item \da{må} ... \da{parkere}（表示“不准/禁止”，必须用 \textit{må ikke}，且后接原形不加 \textit{at}）
  \item \da{bør} / \da{må} ... \da{tage}（接 \textit{hellere} 时两词都地道，常说 \textit{må hellere tage} 或 \textit{bør tage}）
  \item \da{kan} ... \da{tale}
  \item \da{Vil} ... \da{køre}（问对方个人偏好与意愿用 \textit{Vil du...?}）
  \item \da{skal} ... \da{møde}
\end{enumerate}

\textbf{练习 B 答案}
\begin{enumerate}
  \item \da{på}（\textit{på cykel} 骑车）
  \item \da{i}（\textit{i bil} 坐/开汽车）
  \item \da{toget}（动词 \textit{tage} 直接带具体交通工具时要加定冠词后缀！）
  \item \da{med} ... \da{med}（介词 \textit{med} 后接零冠词公共交通）
  \item \da{til}（\textit{til fods} 步行）
\end{enumerate}

\textbf{练习 C 答案}
\begin{enumerate}
  \item \textit{I dag \dabold{kan vi ikke} tage på arbejde på cykel.}（\textit{I dag} 放句首引发 V2 倒装，\textit{ikke} 夹在 \textit{vi} 和 \textit{tage} 之间）
  \item \textit{Med metroen \dabold{skal vi ikke} skifte.} / \textit{Vi skal ikke skifte med metroen.}
  \item \textit{I myldretiden \dabold{bør man ikke} køre i bil til centrum.}
  \item \textit{Vil du hellere køre med tog eller med metro?}
  \item \textit{Man \dabold{må ikke} parkere gratis her, så mange tager bussen.}
\end{enumerate}

\textbf{Før du læser 参考回答 \& 口语复述范文（Resumé）}
\begin{itemize}
  \item \textbf{问题 1}：\textit{De vil ikke cykle, fordi det regner kraftigt, og vinden blæser. Og de bør ikke tage bussen i myldretiden, fordi den ofte sidder fast i trafikken.}
  \item \textbf{问题 2}：\textit{De vælger at tage metroen til arbejde. De kan gå til fods hen til metrostationen med en paraply på otte minutter.}
  \item \textbf{口语复述范文}：\textit{Det er en regnfuld morgen, og vinden blæser kraftigt. Emma vil ellers helst cykle til arbejde for at få motion, men de kan ikke cykle i det dårlige vejr. De bør heller ikke køre med bus i myldretiden, fordi trafikken står stille. Derfor bliver de enige om at tage metroen, hvor de ikke skal skifte. Heldigvis tager det kun otte minutter at gå til fods hen til stationen.}
\end{itemize}
\end{answerbox}
